\section{引言}

纳米硒（selenium nanoparticles，SeNPs）的细胞毒性评价需要同时考虑材料性质与暴露条件，单个粒径或单一浓度下的实验结果难以覆盖不同使用情境。已有细胞与动物研究显示，SeNPs 的生物学效应并不能仅由“纳米硒”这一材料名称概括，材料修饰与暴露水平都需要纳入安全性评价\cite{wang2020prooxidation,lunkov2023chitosan}。对不同表面修饰材料获取剂量—时间响应，有助于明确观测差异出现在哪些条件下，也为预测模型提供直接的评价对象。

文献数据为计算毒理学提供了覆盖多种材料和细胞体系的研究基础。Gul 等整合纳米颗粒材料、细胞和实验属性，通过关联规则分析研究细胞毒性之间的关系 \cite{gul2021association}；Masarkar 等进一步利用随机森林开展跨材料的细胞存活率回归与毒性分类 \cite{masarkar2024}。这些工作说明表格化理化性质和实验条件可用于计算分析。但通用纳米材料上的结果不能直接说明模型对 SeNPs 或新的表面修饰材料具有同等精度，数据来源与验证问题需要分别界定。

历史毒性数据通常在一篇文献内包含同种材料的多个浓度和时间条件，随机划分记录可能把相近实验条件同时分配到训练与验证部分。特征筛选和调参还可能进一步影响性能估计，选择模型的验证结果与报告模型性能的结果应当分离 \cite{cawley2010selection}。对于 SeNPs，另一个问题是输入信息的可获得性：粒径、Zeta 电位与暴露条件通常早于完整毒性实验获得，实测 IC50 却由毒性响应决定，把它作为输入会改变模型的实际使用前提。化学表征带来的预测增量应在相同数据、分组与评价口径下检验，不能只依据训练集的高拟合度判断。

本研究围绕两个相互衔接的问题展开：在按文献分组的历史SeNPs数据上，基础模型能达到什么预测水平，加入化学信息、剂量—时间交互或通用数据迁移能否稳定降低误差；在独立的HUVEC实验中，9种SeNPs的存活率如何随浓度和时间变化，不同指标如何反映材料间差异。研究采用通用基准、历史分组建模、独立HUVEC观测和解盲前封存预测四部分证据，主模型不使用实测IC50，新增实验不并入历史训练集。研究目标是厘清现有信息对预测与实验解释的支持范围，而新材料的严格前瞻预测评价仍以预测值与实测值的独立配对为前提。
